\documentclass[]{../../../../tex/classes/styledArticle}
\usepackage{../../../../tex/packages/hebrewSupport}


\author{שחר פרץ}
\title{\textit{מטלה לסיום שיעור 4} $\sim$ חברה דמוקרטית, תרבות וחינוך}
\begin{document}
	\maketitle
	
	ראשית, סליחה על ההגשה המאוחרת – אני רגיל לתרגילים של הפקולטה מדויקים שניתנים להגשה עד 23:59. 
	
	
	\textbf{משימה: }בחרו אחד מהראיונות שצפיתם בהם והעלו – נקודה שסקרנה אתכם, השאירה עליכם חותם וכתבו עליה כמה מילים. 
	
	\textbf{הדיון: }שיחה עם פרופ' חני לרנר – ההיבטים המוסדיים של דמוקרטיה ותפקידן של חוקות
	
	\textbf{תשובה: }ארצה לדבר על הנושא של חוקה כתובה לעומת חוקה לא כתובה. בסרטון הועלה הנושא בדבר חוקה כתובה, ופרופ' לרנר ציינה שקיימת כזו בכמעט כל המדינות (החברות באו''ם) למעט ישראל, בריטניה וניו־זילנד. לצד זאת היא אמרה שחוקה לא כתובה (שנהוגה באמצעות תקדימים משפטיים, נהגים של מערכות המדינה, וכו') נהוגה במידה כזו או אחרת גם במדינות בהן ישנה חוקה כתובה. 
	
	לאורך התשובה הזו אגדיר חוקה כמכלול פעולות שמתווה את פעילותה של המדינה (הגדרה שאני המצאתי בזה הרגע), שכן בראיון עצמו לא מסופקת הגדרה של חוקה בכלליות, אלא רק של חוקה כתובה ושל חוקה שאינה כתובה. 
	
	הסיבה שבחרתי להתייחס לנקודה זו, היא בגלל ההתנגשות בשנים האחרונות בין הקואליציה לבין החוקה הלא־כתובה (או החצי־כתובה, בדמות חוקי יסוד, אך אלו בעלי מעמד משפטי כמעט זהה לחוקים רגילים: מתקבלים באותו הרוב, וניתנים לביטול ע''י בית המשפט). מצד אחד, החוקה הנהוגה בישראל אינה כתובה שחור על גבי לבן;
	\begin{itemize}
		\item כך גם הנהגים הרבים של בית המשפט (לדוגמה, מינוי נשיא העליון להיות האדם הוותיק ביותר, תקדימים משפטיים וכו') אינם כתובים שחור־על גבי לבן באף מסמך בעל מעמד חוקתי. 
		\item גם רבים מהנהגים של מערכות המדינה (לדוגמה, שמעורבות השר לבטחון לאומי מוגבלת לדרגת בכירות מסוימת, מתחתיה מניחים למפכ''ל לפעול לפי שיקול לדעתו). 
	\end{itemize}
	מצד שני, נקודה שהולכת ונשנת בדיון על החוקה (הכתובה והלא כתובה) הוא תפקידה להגביל את כוח מוסדות השלטון, דהיינו התעלמות מהחוקה הלא־כתובה של מדינת ישראל חסרת החוקה הכתובה תגרור התמוססות של כל הבלמים המגבילים את הרוב בכנסת. 
	
	הסיבה שאני חושב שהנקודה הזו כל־כך חשובה, היא שלדעתי שני הצדדים אינם מבינים אותה כראוי: 
	\begin{itemize}
		\item מצד אחד, אלו שגורסים שיש להתעלם מאותם הנהגים שכן אינם מצויים בחוק, או שאינם מבינים את תפקידה של החוקה במדינה (חושבים שדמוקרטיה נגמרת בשלטון מוחלט, או יש שיגידו עריצות, של הרוב), או שאינם מבינים את התפקיד שממלאת החוקה הבלתי־כתובה במדינת ישראל. 
		\item מאידך, אלו שמתנגדים לפעולות הממשלה להרחבת סמכויתיה בצורה המאפשרת את עריצות הרוב, לטענתי לא מבינים מה מקור הבעיה. רבים יצביעו על הכשלים ברפורמה המשפטית, או על נסיונתם של חברי הכנסת להתערב בסמכויתהן של רשויות עצמאיות, אך במצב היפוטתי בו האופוזציה נבחרת, מה הלאה? מה מקור הבעיה? מקור (לפחות חלקי) לבעיה הוא ערעור על החוקה הלא כתובה, שמתאפשר מעצם היותה בלתי כתובה. אז, הכיצד ש(ככל הידוע לי) אף אדם לא הציע להפוך, לפחות חלקים (כמו נהגי בית המשפט) לחוק כתוב? 
	\end{itemize}
	
	לסיכום, אני חושב ששני הצדדים (או לכל הפחות, האדם הממוצע בכל צד בנפרד) של המפה הפוליטית לא מבינים את חשיבותה של החוקה בלתי כתובה. צד אחד טוען שאין לממש אותה שכן איננה חוק, וצד אחר לא מודע מספיק לבלמים שהיא מספקת כנגד אותן הפעולות שהוא מבקר. \hfill $\blacksquare$
	
	\textit{הערה: }יכול להיות שיש טעויות עובדתיות בטקסט לעיל. אני כותב על בסיס הידע הכללי שלי ולא בדקתי כל עובדה בנפרד. אם מצאתם כאלו, אשמח שתתקנו אותי. 
	
	
	\ndoc
\end{document}